\section*{Hoofdstuk \ref{ch:b1:gene-expression} --- Genexpressie: transcriptie en translatie}

\begin{solution}{exo:b1:gene-expression:1}
\omterm{def:b1:nucleic-acids:chain}{RNA} $5'$-AUGCCUAAG-$3'$; coderende streng $5'$-ATGCCTAAG-$3'$.
\end{solution}

\begin{solution}{exo:b1:gene-expression:2}
Een cap aan het $5'$-uiteinde (een gewijzigde guanine), een \omterm{prop:b1:gene-expression:processing}{poly-A-staart} aan het
$3'$-uiteinde, en het verwijderen van de \omterm{def:b1:genomes:gene}{intronen} door \omterm{prop:b1:gene-expression:processing}{splicing}.
\end{solution}

\begin{solution}{exo:b1:gene-expression:3}
AUG CCG AAA GUU UGA: Met-Pro-Lys-Val, daarna stop.
\end{solution}

\begin{solution}{exo:b1:gene-expression:4}
A: het binnenkomende geladen tRNA wordt toegelaten en gecontroleerd. P: het tRNA
dat de groeiende keten draagt; de binding ontstaat tussen zijn keten en het
\omterm{def:b1:proteins:aminoacid}{aminozuur} in de A-plaats. E: het geleegde tRNA op weg naar buiten.
\end{solution}

\begin{solution}{exo:b1:gene-expression:5}
Coderend $2400 - 600 = 1800$ \omterm{def:b1:nucleic-acids:nucleotide}{nucleotiden}: 599 residuen (600 \omterm{def:b1:gene-expression:code}{codons}, waarvan één een
stop). Tijd $599/4 = \qty{150}{s}$. \omterm{def:b1:gene-expression:ribosome}{Ribosomen} $1800/90 = 20$ tegelijk.
\end{solution}

\begin{solution}{exo:b1:gene-expression:6}
De boodschap wordt vanaf een vast startpunt in drietallen gelezen; één of twee
extra basen verschuiven elk \omterm{def:b1:gene-expression:code}{codon} erna en de rest van het \omterm{def:b1:proteins:peptide}{eiwit} is onzin, meestal
eindigend bij een voortijdige stop. Drie extra basen voegen één \omterm{def:b1:gene-expression:code}{codon} toe en
herstellen het raam: het \omterm{def:b1:proteins:peptide}{eiwit} heeft één residu extra en enkele verkeerde tussen de
invoegingen, en werkt vaak nog.
\end{solution}

\begin{solution}{exo:b1:gene-expression:7}
CAG (Gln) naar UAG (stop): de keten eindigt bij residu 150, een afgekapt en vrijwel
zeker inactief \omterm{def:b1:proteins:peptide}{eiwit} (nonsensmutatie). CAG naar CAA: nog steeds Gln, stil. CAG naar
CGG: Arg in plaats van Gln, een vervanging (missense) waarvan het effect van de
plaats afhangt.
\end{solution}

\begin{solution}{exo:b1:gene-expression:8}
Het \omterm{def:b1:gene-expression:ribosome}{ribosoom} controleert alleen de paring van \omterm{def:b1:gene-expression:code}{codon} en \omterm{def:b1:gene-expression:trna}{anticodon}; het kan het
\omterm{def:b1:proteins:aminoacid}{aminozuur} niet zien. Wordt cysteïne op haar tRNA langs chemische weg in alanine
omgezet, dan bouwt het \omterm{def:b1:gene-expression:ribosome}{ribosoom} alanine in overal waar de boodschap cysteïne zegt:
het tRNA wordt gelezen, niet het \omterm{def:b1:proteins:aminoacid}{aminozuur}. Dus is het synthetase, dat elk \omterm{def:b1:proteins:aminoacid}{aminozuur}
met het juiste tRNA koppelt, de eigenlijke vertaler.
\end{solution}

\begin{solution}{exo:b1:gene-expression:9}
$400\times 4 = 1600$ bindingen; de twee GTP per residu van het \omterm{def:b1:gene-expression:ribosome}{ribosoom} zijn de
helft, de twee van de synthetasen de andere helft.
\end{solution}

\begin{solution}{exo:b1:gene-expression:10}
Bij eukaryoten wordt het transcript in de celkern gemaakt en zitten de \omterm{def:b1:gene-expression:ribosome}{ribosomen} in
het \omterm{def:b1:eukaryotic-cell:organelle}{cytosol}, gescheiden door de envelop; en het transcript is pas een boodschapper
nadat het is gespliced en van een cap voorzien. Die scheiding maakt de bewerking
zelf mogelijk --- alternatieve \omterm{prop:b1:gene-expression:processing}{splicing}, de sturing van de uitvoer, en een controle
dat de boodschap volledig is voordat zij wordt gelezen.
\end{solution}

\begin{solution}{exo:b1:gene-expression:11}
Vier boodschappers: met beide \omterm{def:b1:genomes:gene}{exonen} (190 \omterm{def:b1:nucleic-acids:nucleotide}{nucleotiden} erbij), met alleen 3 (90),
met alleen 4 (100), met geen van beide (0). Het raam blijft behouden als de
toegevoegde lengte een veelvoud van 3 is: 0 en 90 behouden het; 100 en 190
verschuiven het, zodat de boodschappers met alleen \omterm{def:b1:genomes:gene}{exon} 4 of met beide \omterm{def:b1:genomes:gene}{exonen} \omterm{def:b1:genomes:gene}{exon}
5 buiten het raam lezen en het \omterm{def:b1:proteins:peptide}{eiwit} afkappen. Alternatieve \omterm{prop:b1:gene-expression:processing}{splicing} moet het raam
eerbiedigen, en juist daarom zijn \omterm{def:b1:genomes:gene}{exonen} vaak een veelvoud van drie.
\end{solution}

\begin{solution}{exo:b1:gene-expression:12}
Willekeurig: niets in de chemie zegt dat UUU Phe moet betekenen; de toewijzingen
zijn afspraken die door de synthetasen worden vastgehouden. Geoptimaliseerd: de
derde plaats is het sterkst gedegenereerd, zodat de fouten en mutaties die daar
vallen vaak stil zijn, en \omterm{def:b1:gene-expression:code}{codons} die in één base verschillen coderen doorgaans voor
verwante \omterm{def:b1:proteins:aminoacid}{aminozuren}, zodat een vervanging vaak mild is --- de code beperkt de schade
van een fout. Bevroren: een synthetase dat zijn specificiteit zou veranderen, zou
in hetzelfde ogenblik elk \omterm{def:b1:proteins:peptide}{eiwit} van de \omterm{def:b1:cell-unit-of-life:cell}{cel} veranderen, duizenden tegelijk; vrijwel
geen enkele \omterm{def:b1:cell-unit-of-life:cell}{cel} overleeft dat, zodat de code niet kan driften en sinds de
gemeenschappelijke voorouder van alle levende wezens vastligt.
\end{solution}

\begin{solution}{pb:b1:gene-expression:1}
\textbf{1.} $30\,000/30 = \qty{1000}{s}$, ongeveer \qty{17}{min}.
\textbf{2.} $28\,000/30\,000 = \qty{93}{\%}$.
\textbf{3.} $30\,000/2000 = 15$.
\textbf{4.} $2\times 30\,000 = \num{60000}$ \omterm{def:b1:water-small-molecules:atp}{ATP}.
\textbf{5.} $1000/10 = 100$ polymerasen op het \omterm{def:b1:genomes:gene}{gen}; 360 boodschappers per uur.
\textbf{6.} Zeven \omterm{def:b1:genomes:gene}{intronen}, gemiddeld $28\,000/7 = \qty{4000}{}$ \omterm{def:b1:nucleic-acids:nucleotide}{nucleotiden}.
\textbf{7.} De \omterm{prop:b1:gene-expression:dogma}{transcriptie} (\qty{17}{min}); de \omterm{def:b1:genomes:gene}{intronen} worden gespliced zodra zij
zijn gemaakt, en het laatste voegt maar een minuut toe.
\textbf{8.} \omterm{def:b1:genomes:gene}{Gen} $30\,000\times 0.34\,\mathrm{nm} = \qty{10}{\micro m}$; boodschapper \qty{0.68}{\micro m}.
\textbf{9.} $500/4 = \qty{125}{s}$.
\textbf{10.} $2000/80 = 25$.
\textbf{11.} Er komt elke $80/4 = \qty{20}{s}$ een keten gereed: 180 per uur.
\textbf{12.} $2/\ln 2 = \qty{2.9}{h}$ volle productie: ongeveer 520 \omterm{def:b1:proteins:peptide}{eiwitten}.
\textbf{13.} $4\times 500 = 2000$ \omterm{def:b1:water-small-molecules:atp}{ATP}.
\textbf{14.} $60\,000/520 = 115$ \omterm{def:b1:water-small-molecules:atp}{ATP} per \omterm{def:b1:proteins:peptide}{eiwit}.
\textbf{15.} Ongeveer 2100 \omterm{def:b1:water-small-molecules:atp}{ATP}, waarvan \qty{95}{\%} in de \omterm{prop:b1:gene-expression:dogma}{translatie}.
\textbf{16.} $1 - (1 - 10^{-5})^{1503} \approx 1503\times 10^{-5} =
\qty{1.5}{\%}$.
\textbf{17.} $1 - (1 - 10^{-4})^{500} \approx \qty{5}{\%}$.
\textbf{18.} \omterm{prop:b1:gene-expression:dogma}{Transcriptie}: $1.5\%\times 0.75\times 0.67 \approx
0.75\%$ van de boodschappers is defect, en dus ook van de \omterm{def:b1:proteins:peptide}{eiwitten}; \omterm{prop:b1:gene-expression:dogma}{translatie}:
$5\%\times 0.67 = 3.3\%$; samen ongeveer \qty{4}{\%} van de moleculen.
\textbf{19.} Een fout in een \omterm{def:b1:proteins:peptide}{eiwit} blijft beperkt tot één molecuul, dat spoedig
wordt afgebroken; een fout in een boodschapper wordt in honderden \omterm{def:b1:proteins:peptide}{eiwitten}
gekopieerd; een fout in de replicatie wordt door elke nakomeling voorgoed geërfd.
De kosten van een fout, en daarmee de nauwkeurigheid die de moeite waard is,
stijgen bij elke stap terug.
\textbf{20.} $3000/180 = 17$ boodschappers aanwezig.
\textbf{21.} $17\times 25 = 420$ \omterm{def:b1:gene-expression:ribosome}{ribosomen}.
\textbf{22.} $\num{3e9}\times 4 = \num{1.2e10}$ \omterm{def:b1:water-small-molecules:atp}{ATP} per dag
$= \qty{2e-14}{mol}$, \qty{1e-9}{J} per dag, oftewel \qty{1.2e-14}{W} ---
\qty{6}{fW} per kubieke micrometer.
\textbf{23.} $10^9$ \omterm{def:b1:water-small-molecules:atp}{ATP} per seconde is $\num{8.6e13}$ per dag: de eiwitsynthese is
daar een honderdste van in deze traag vernieuwende \omterm{def:b1:cell-unit-of-life:cell}{cel} (in een delende bacterie de
helft).
\textbf{24.} Geen rijpe boodschapper: de primaire transcripten hopen zich in de
celkern op en het \omterm{def:b1:proteins:peptide}{eiwit} verdwijnt binnen enkele uren naarmate zijn boodschappers
vervallen. Het bacteriële \omterm{def:b1:proteins:peptide}{eiwit}, waarvan het \omterm{def:b1:genomes:gene}{gen} geen \omterm{def:b1:genomes:gene}{intronen} heeft, blijft
onaangetast.
\textbf{25.} Ongeveer 2100 \omterm{def:b1:water-small-molecules:atp}{ATP} per molecuul --- 2000 voor de \omterm{prop:b1:gene-expression:dogma}{translatie}, honderd
voor zijn aandeel in het transcript; het eerste \omterm{def:b1:proteins:peptide}{eiwit} na \qty{1000}{s} \omterm{prop:b1:gene-expression:dogma}{transcriptie},
een minuut bewerking en uitvoer, en \qty{125}{s} \omterm{prop:b1:gene-expression:dogma}{translatie}: ongeveer twintig
minuten.
\end{solution}
